\cvsection{Extracurriculars}
\begin{cventries}
  \cventry
    {Event Coordinator, Technology Lead}
    {Solastalgia}
    {Vancouver, BC, Canada}
    {July 2023 - Present}
    {
      \begin{cvitems}
        \item Joined a youth-led organization to develop a community arts zine exploring eco-emotions, resulting in a successful launch supported by OceanWise Canada
        \item Led development of a new Nature walk event, in collaboration with Pacific Spirit Park Society, with an attendance of 20+ community members
        \item Established a partnership with Moberly Arts \& Cultural centre as a space to host events and meetings, with 10+ hours of community engagement
      \end{cvitems}
    }
  \cventry
    {RISE Ambassador Vancouver}
    {Apathy is Boring}
    {Vancouver, BC, Canada}
    {Jan. 2023 - June 2023}
    {
      \begin{cvitems}
        \item Worked among a team of 11 youth ambassadors to consult community needs and design an youth-led event tackling the theme of "Rsilient Cities".
        \item Delivered the final project as a community crawl with youth ages 18-30 around Commercial Drive, highlighting historical events and injustices that shaped the neighbourhood.
        \item Designed an interactive map of the route taken that serves as a future
        resource for youth.
      \end{cvitems}
    }
  \cventry
    {Youth Climate Innovation Lab - Cohort 1 Participant}
    {CityHive Vancouver}
    {Vancouver, BC, Canada}
    {Sep 2022 - Dec 2022}
    {
      \begin{cvitems}
        \item Led an EnROADS workshop for the second YCIL cohort, educating
        participants about climate policies and actions which make a real difference
        \item Attended educational and activating workshops to develop my understanding of climate action within the Vancouver context along with
        20+ engaged community youth members.
        \item Implemented a project with the Vancouver Economic Commission to explore how hardware stores could become "Climate Hardware Shops".
        \item Presented final project findings at the community event to a virtual audience of 50+ community collaborators and organization.
      \end{cvitems}
    }
  % \cventry
  %   {Physics PASS Leader}
  %   {UBC Peer Assisted Study Sessions (PASS)}
  %   {Vancouver, BC, Canada}
  %   {Sept. 2014 - Apr. 2015}
  %   {
  %     \begin{cvitems}
  %       \item Adapted strong inter-personal skills to design intimate, interactive study sessions for Physics with first year students every week.
  %       \item Administered the first-ever PASS program Facebook group to fill in a communication gap between students and
  %       peer leaders.
  %       \item Independently planned and presented a leadership workshop as the sole PASS representative at the UBC Student-Leadership Conference, one of Canada’s largest student-run conferences.
  %     \end{cvitems}
  %   }
  % \cventry
  %   {Student Leader - Academic Coach}
  %   {UBC Jumpstart}
  %   {Vancouver, BC, Canada}
  %   {Aug. 2013 - Apr. 2014}
  %   {
  %     \begin{cvitems}
  %       \item Collaborated with a team of 50+ other student leaders during a two week immersion program for 1000+ international students.
  %       \item Worked closely with a Faculty Fellow to coach a personal group of 30+ students in a Learning Community, employing active listening and inter-personal skills while further developing networking and teamwork skills.
  %       \item Planned and coordinated a professional advice seminar for 50 students, in collaboration with faculty, to provide students a chance to think about careers and visit work-places.
  %       \item Spearheaded the creation of a secret lip-dub video for a leaving staff member, involving a large group of students leaders and staff members.
  %     \end{cvitems}
  %   }
  % \cventry
  %   {Associate Curriculum Coordinator}
  %   {UBC Engineering Undergraduate Society}
  %   {Vancouver, BC, Canada}
  %   {Apr. 2013 - Apr. 2014}
  %   {
  %     \begin{cvitems}
  %       \item Selected to be involved in the refinement of the first and second year curriculum, establishing direct improvements and effects for the student body, resulting in the reception of the One Year Service with Gratitude Award.
  %       \item Created Outgoing First-Year Survey to assess the student body’s needs, gathering information from 87 students to act as feedback for curriculum changes.
  %       \item Summarized survey results into a 21 page Report, highlighting general problems with the First-Year curriculum and suggested several realistic, course-specific solutions.
  %       \item Collaborated with professors to create a new set of surveys for each Case Study in APSC 150, a first year engineering course.
  %     \end{cvitems}
  %   }
\end{cventries}
